\documentclass[11pt,a4paper]{article}

% --- Paquetes ---
\usepackage[utf8]{inputenc}
\usepackage[T1]{fontenc}
\usepackage[spanish]{babel}
\usepackage{amsmath, amssymb, amsthm}
\usepackage{geometry}
\usepackage{tikz}
\usepackage{pgfplots}
\usepackage{enumitem}
\usepackage{hyperref}
\usepackage{xcolor}
\usepackage{tcolorbox}

\geometry{margin=2.5cm}
\pgfplotsset{compat=1.18}

% --- Definiciones ---
\newcommand{\barvar}[1]{\overline{#1}}

\newtcolorbox{definicionbox}[1][]{colback=yellow!10, colframe=red!70!black, title=#1, fonttitle=\bfseries}
\newtcolorbox{resultadobox}[1][]{colback=yellow!10, colframe=red!80!black, title=#1, fonttitle=\bfseries}
\newtcolorbox{explicacionbox}[1][]{colback=red!5, colframe=red!60!black, title=#1, fonttitle=\bfseries}
\newtcolorbox{puntosclavebox}[1][]{colback=yellow!10, colframe=green!50!black, title=#1, fonttitle=\bfseries}

% --- Documento ---
\title{\textbf{Clase 3: Multiplicador del Gasto, Política Fiscal y Curva IS} \\ \large Macroeconomía II}
\author{Universidad Católica del Norte}
\date{29 de agosto, 2026}

\begin{document}
\maketitle
\tableofcontents
\newpage

% ============================================================
\section{El Multiplicador del Gasto}
% ============================================================

\begin{definicionbox}[Pregunta central]
¿En cuánto aumenta el nivel de equilibrio de producción cuando existe un cambio exógeno en el gasto autónomo de una unidad monetaria?
\end{definicionbox}

\subsection{Derivación del Multiplicador}

Consideremos las variaciones sucesivas en $Y$ cuando $\barvar{A}_0$ cambia:

\begin{center}
\begin{tabular}{lll}
\textbf{Variación en $\Delta \barvar{A}_0$} & & \textbf{Variación en $\Delta Y$} \\
\hline
1\textsuperscript{a} Variación & $\rightarrow$ & $\Delta \barvar{A}_0$ \quad (cambio exógeno) \\
2\textsuperscript{a} Variación & $\rightarrow$ & $c \cdot \Delta \barvar{A}_0$ \\
3\textsuperscript{a} Variación & $\rightarrow$ & $c \cdot (c \cdot \Delta \barvar{A}_0) = c^2 \cdot \Delta \barvar{A}_0$ \\
4\textsuperscript{a} Variación & $\rightarrow$ & $c^3 \cdot \Delta \barvar{A}_0$ \\
\vdots & & \vdots \\
\end{tabular}
\end{center}

¿Qué ocurre cuando aumenta $\barvar{A}_0$?
\begin{align}
\uparrow \barvar{A}_0 \;\Rightarrow\; &\uparrow DA \;(\Delta \barvar{A}_0) \;\Rightarrow\; \uparrow Y \;(\Delta \barvar{A}_0) \notag \\
\Rightarrow\; &\uparrow Y_d \;(c\,\Delta \barvar{A}_0) \;\Rightarrow\; \uparrow C \;(c\,\Delta \barvar{A}_0) \notag \\
\Rightarrow\; &\uparrow DA \;(c\,\Delta \barvar{A}_0) \;\Rightarrow\; \uparrow Y \;(c\,\Delta \barvar{A}_0) \notag \\
&\quad \cdots \notag
\end{align}

La variación total en $Y$ es la suma de todas las variaciones:

\begin{resultadobox}
\[
\Delta Y = \Delta \barvar{A}_0 + c\,\Delta \barvar{A}_0 + c^2\,\Delta \barvar{A}_0 + c^3\,\Delta \barvar{A}_0 + \cdots
\]
\end{resultadobox}

Factorizamos $\Delta \barvar{A}_0$:
\[
\Delta Y = \Delta \barvar{A}_0 \left(1 + c + c^2 + c^3 + c^4 + \cdots \right)
\]

Dividiendo por $\Delta \barvar{A}_0$:
\begin{equation}
\frac{\Delta Y}{\Delta \barvar{A}_0} = 1 + c + c^2 + c^3 + c^4 + \cdots \tag{1}
\end{equation}

\begin{itemize}
    \item El lado izquierdo representa: ¿cuánto varía $\Delta Y$ cuando $\Delta \barvar{A}_0$ varía en una unidad?
    \item El lado derecho es una \textbf{suma infinita} (que se va haciendo cada vez más pequeña, ya que $0 < c < 1$).
\end{itemize}

\subsubsection{Resolución de la serie infinita}

Multiplicamos la ecuación (1) por $(-c)$:
\begin{equation}
-c \cdot \frac{\Delta Y}{\Delta \barvar{A}_0} = -c - c^2 - c^3 - c^4 - c^5 - \cdots \tag{2}
\end{equation}

Sumamos $(1) + (2)$:
\[
\frac{\Delta Y}{\Delta \barvar{A}_0} - c \cdot \frac{\Delta Y}{\Delta \barvar{A}_0} = 1
\]

Factorizamos:
\[
(1 - c) \cdot \frac{\Delta Y}{\Delta \barvar{A}_0} = 1
\]

\begin{resultadobox}[Multiplicador del Gasto]
\begin{equation}
\frac{\Delta Y}{\Delta \barvar{A}_0} = \frac{1}{1-c} = \alpha_m \tag{3}
\end{equation}
donde $\alpha_m$ es el \textbf{Multiplicador del Gasto}.
\end{resultadobox}

\subsection{Interpretación del Multiplicador}

\begin{definicionbox}[Definición del Multiplicador]
El multiplicador ($\alpha_m$) es el monto en el cual cambia la producción de equilibrio cuando la demanda agregada autónoma aumenta en 1 unidad.
\[
\alpha_m = \frac{1}{1-c}
\]
\end{definicionbox}

Dado que $0 < c < 1$:
\[
\frac{1}{1-c} > 1
\]

Es decir, el multiplicador es siempre \textbf{mayor que 1}. Si el gasto autónomo varía en $\Delta \barvar{A}_0$, la producción de equilibrio varía en:

\begin{resultadobox}
\[
\Delta Y = \alpha_m \cdot \Delta \barvar{A}_0 = \frac{1}{1-c} \cdot \Delta \barvar{A}_0 \quad (> \Delta \barvar{A}_0)
\]
\end{resultadobox}

\subsection{El Multiplicador en el gráfico}

\begin{center}
\begin{tikzpicture}
    \begin{axis}[
        axis lines=middle,
        xlabel={$Y$}, ylabel={$DA$},
        xmin=-1, xmax=12, ymin=-1, ymax=12,
        xtick=\empty, ytick=\empty,
        every axis x label/.style={at={(ticklabel* cs:1)}, anchor=west},
        every axis y label/.style={at={(ticklabel* cs:1)}, anchor=south},
        width=11cm, height=9cm,
    ]
    % Línea de 45 grados
    \addplot[domain=0:11, thick, blue] {x};
    \node[above left, blue] at (axis cs:10.5,10.5) {$DA = Y$};

    % DA con Ao1
    \addplot[domain=0:11, thick, red!70!black] {2 + 0.6*x};

    % DA con Ao2
    \addplot[domain=0:11, thick, red!70!black] {4 + 0.6*x};

    % Equilibrios
    \addplot[only marks, mark=*, mark size=3pt, teal] coordinates {(5, 5)};
    \node[above left, teal] at (axis cs:5,5) {$E_1$};

    \addplot[only marks, mark=*, mark size=3pt, teal] coordinates {(10, 10)};
    \node[above left, teal] at (axis cs:10,10) {$E_2$};

    % Líneas punteadas
    \addplot[dashed, gray] coordinates {(5,0) (5,5)};
    \addplot[dashed, gray] coordinates {(10,0) (10,10)};

    % Delta Ao
    \draw[->, thick, green!50!black] (axis cs:9.5,9.7) -- (axis cs:9.5,7.7);
    \node[right, green!50!black] at (axis cs:9.6,8.7) {$\Delta \barvar{A}_0$};

    % Delta Y
    \draw[<->, thick, orange] (axis cs:5,-0.5) -- (axis cs:10,-0.5);
    \node[below, orange] at (axis cs:7.5,-0.5) {$\Delta Y$};

    % Etiquetas
    \node[left] at (axis cs:0,2) {$\barvar{A}_0^1$};
    \node[left] at (axis cs:0,4) {$\barvar{A}_0^2$};
    \node[below] at (axis cs:5,0) {$Y_1$};
    \node[below] at (axis cs:10,0) {$Y_2$};

    \end{axis}
\end{tikzpicture}
\end{center}

Si por alguna razón aumenta el gasto autónomo (por ejemplo $\uparrow \barvar{I}$), lo que produce un aumento en el ingreso, la gente gastará más, lo que aumentará aún más el ingreso vía el consumo.

\begin{puntosclavebox}[Puntos Clave]
\begin{itemize}
    \item $\uparrow \barvar{A}_0 \;\Rightarrow\; \uparrow Y^*$
    \item $\dfrac{\Delta Y}{\Delta \barvar{A}_0} = \alpha_m$
    \item Entre mayor sea $c \;\Rightarrow\; \uparrow \alpha_m$
\end{itemize}
\end{puntosclavebox}

% ============================================================
\section{Determinación de la Producción de Equilibrio}
% ============================================================

¿Cómo alcanza la economía el equilibrio?

\subsection{Caso 1: Exceso de Oferta ($Y > DA$)}

\begin{center}
\begin{tikzpicture}
    \begin{axis}[
        axis lines=middle,
        xlabel={$Y$}, ylabel={$DA$},
        xmin=-1, xmax=12, ymin=-1, ymax=12,
        xtick=\empty, ytick=\empty,
        every axis x label/.style={at={(ticklabel* cs:1)}, anchor=west},
        every axis y label/.style={at={(ticklabel* cs:1)}, anchor=south},
        width=11cm, height=8cm,
    ]
    % Línea de 45 grados
    \addplot[domain=0:11, thick, blue] {x};
    \node[above left, blue] at (axis cs:10,10) {$DA = Y$};

    % DA
    \addplot[domain=0:11, thick, red!70!black] {2 + 0.6*x};
    \node[right, red!70!black] at (axis cs:9,7.4) {$DA = \barvar{A}_0 + cY$};

    % Equilibrio
    \addplot[only marks, mark=*, mark size=3pt, orange] coordinates {(5, 5)};
    \node[above left, orange] at (axis cs:5,5) {$E$};
    \node[above, orange] at (axis cs:5,5.5) {\small Equilibrio Estable};

    % Punto Y1 > Y*
    \addplot[only marks, mark=*, mark size=3pt, black] coordinates {(8, 8)};
    \addplot[only marks, mark=*, mark size=3pt, black] coordinates {(8, 6.8)};

    % Líneas punteadas Y1
    \addplot[dashed, gray] coordinates {(8,0) (8,8)};

    % Flechas de ajuste hacia equilibrio
    \draw[->, thick, teal] (axis cs:7.8,7.5) -- (axis cs:6.5,6.5);
    \draw[->, thick, teal] (axis cs:6.3,6.3) -- (axis cs:5.5,5.5);

    % Etiquetas
    \node[left] at (axis cs:0,2) {$\barvar{A}_0$};
    \node[below] at (axis cs:5,0) {$Y^*$};
    \node[below] at (axis cs:8,0) {$Y_1$};
    \node[left] at (axis cs:7.3,5.5) {$\longleftarrow$};

    \end{axis}
\end{tikzpicture}
\end{center}

En el punto $Y_1$ no estamos en equilibrio. La producción es mayor que la demanda agregada planeada: \textbf{Exceso de Oferta}.

$Y > DA \;\Rightarrow$ Las empresas producen más de lo que la economía demanda:

\begin{enumerate}
    \item \textbf{Muy corto plazo:} $\uparrow$ Inventario ($IU > 0$)
    \item $\Downarrow$ $\downarrow$ Contratación ($\downarrow L$)
    \item $\Downarrow$ $\downarrow$ Ingreso ($\downarrow Y$)
    \item $\Downarrow$ $\downarrow$ Demanda ($\downarrow DA$)
    \item $\Downarrow$ $\downarrow Y \;\Rightarrow\; \downarrow DA \;\cdots$
    \item $\Downarrow$ Ajuste (Equilibrio)
\end{enumerate}

\subsection{Caso 2: Exceso de Demanda ($Y < DA$)}

\begin{center}
\begin{tikzpicture}
    \begin{axis}[
        axis lines=middle,
        xlabel={$Y$}, ylabel={$DA$},
        xmin=-1, xmax=12, ymin=-1, ymax=12,
        xtick=\empty, ytick=\empty,
        every axis x label/.style={at={(ticklabel* cs:1)}, anchor=west},
        every axis y label/.style={at={(ticklabel* cs:1)}, anchor=south},
        width=11cm, height=8cm,
    ]
    % Línea de 45 grados
    \addplot[domain=0:11, thick, blue] {x};
    \node[above left, blue] at (axis cs:10,10) {$DA = Y$};

    % DA
    \addplot[domain=0:11, thick, red!70!black] {2 + 0.6*x};
    \node[right, red!70!black] at (axis cs:9,7.4) {$DA = \barvar{A}_0 + cY$};

    % Equilibrio
    \addplot[only marks, mark=*, mark size=3pt, orange] coordinates {(5, 5)};
    \node[above left, orange] at (axis cs:5,5) {$E$};
    \node[above, orange] at (axis cs:5,5.5) {\small Equilibrio Estable};

    % Punto Y2 < Y*
    \addplot[only marks, mark=*, mark size=3pt, black] coordinates {(2, 2)};
    \addplot[only marks, mark=*, mark size=3pt, black] coordinates {(2, 3.2)};

    % Líneas punteadas Y2
    \addplot[dashed, gray] coordinates {(2,0) (2,3.2)};

    % Flechas de ajuste hacia equilibrio
    \draw[->, thick, orange!80!red] (axis cs:2.5,2.8) -- (axis cs:3.5,3.8);
    \draw[->, thick, orange!80!red] (axis cs:3.7,4.0) -- (axis cs:4.5,4.5);

    % Etiquetas
    \node[left] at (axis cs:0,2) {$\barvar{A}_0$};
    \node[below] at (axis cs:5,0) {$Y^*$};
    \node[below] at (axis cs:2,0) {$Y_2$};

    \end{axis}
\end{tikzpicture}
\end{center}

En el punto $Y_2$ no estamos en equilibrio. La producción es menor que la demanda agregada planeada: \textbf{Exceso de Demanda}.

$Y < DA \;\Rightarrow$ Las empresas producen menos de lo que la economía demanda:

\begin{enumerate}
    \item \textbf{Muy corto plazo:} $\downarrow$ Inventario ($IU < 0$)
    \item $\Downarrow$ $\uparrow$ Contratación ($\uparrow L$)
    \item $\Downarrow$ $\uparrow$ Ingreso ($\uparrow Y$)
    \item $\Downarrow$ $\uparrow$ Demanda ($\uparrow DA$)
    \item $\Downarrow$ $\uparrow Y \;\Rightarrow\; \uparrow DA \;\cdots$
    \item $\Downarrow$ Ajuste (Equilibrio)
\end{enumerate}

% ============================================================
\section{Sector Gubernamental y Política Fiscal}
% ============================================================

\begin{definicionbox}[Preguntas centrales]
¿Cómo puede hacer el gobierno para aumentar la producción de equilibrio? ¿Qué herramientas tiene el gobierno para alterar el crecimiento?
\end{definicionbox}

\textbf{Política Fiscal:} Herramientas que posee el gobierno para cambiar el curso del crecimiento económico en el corto plazo.

Las herramientas de la política fiscal son:

\begin{center}
\begin{tikzpicture}
    \node[draw, thick, red!70!black, fill=red!5, minimum width=5cm, minimum height=1cm] (PF) at (0,0) {\textbf{POLÍTICA FISCAL}};

    \node[below left=1.5cm and 2cm of PF] (G) {$\barvar{G}$};
    \node[below=1.5cm of PF] (Tr) {$\barvar{Tr}$};
    \node[below right=1.5cm and 2cm of PF] (t) {$t \cdot Y$};

    \draw[->, thick, violet] (PF.south west) -- (G);
    \draw[->, thick, violet] (PF.south) -- (Tr);
    \draw[->, thick, violet] (PF.south east) -- (t);

    \node[below=0.2cm of G, align=center] {\small Gasto \\ \small Gubernamental};
    \node[below=0.2cm of Tr, align=center] {\small Transferencias};
    \node[below=0.2cm of t, align=center] {\small Impuesto \\ \small a la renta};
\end{tikzpicture}
\end{center}

\subsection{Impuesto a la Renta: Reescribiendo la función de Consumo}

Ahora introducimos un \textbf{impuesto proporcional a la renta} ($t \cdot Y$) en lugar del impuesto de suma fija ($\barvar{T}$).

La función de consumo es:
\[
C = \barvar{C} + c \, Y_d
\]

El nuevo ingreso disponible es:
\[
Y_d = Y + \barvar{Tr} - t \cdot Y = (1-t)Y + \barvar{Tr}
\]

Sustituyendo:
\begin{align}
C &= \barvar{C} + c\left(Y + \barvar{Tr} - tY\right) \notag \\
C &= \barvar{C} + cY + c\barvar{Tr} - ctY \notag
\end{align}

\begin{resultadobox}[Nueva Función de Consumo]
\begin{equation}
C = \barvar{C} + c\barvar{Tr} + c(1-t)Y \tag{4}
\end{equation}
donde $c(1-t)$ es la fracción del ingreso después de los impuestos que se destina al consumo.
\end{resultadobox}

\subsection{Nueva Función de Demanda Agregada}

¿Cómo afecta esto a la Demanda Agregada? Recordando $DA = C + I + G$:
\begin{align}
DA &= \barvar{C} + c\barvar{Tr} + c(1-t)Y + \barvar{I} + \barvar{G} \notag
\end{align}

Reordenamos:
\[
DA = \underbrace{\barvar{C} + c\barvar{Tr} + \barvar{I} + \barvar{G}}_{\barvar{A}_0 \;\text{(Nuevo Gasto Autónomo)}} + c(1-t)Y
\]

\begin{resultadobox}[Nueva Función de Demanda Agregada]
\begin{equation}
DA = \barvar{A}_0 + c(1-t)Y \tag{5}
\end{equation}
donde:
\begin{itemize}
    \item $\barvar{A}_0 = \barvar{C} + c\barvar{Tr} + \barvar{I} + \barvar{G}$ es el \textbf{Gasto Autónomo}.
    \item $c(1-t)$ es la nueva \textbf{Propensión Marginal a Consumir} (ajustada por impuestos).
    \item $t$ es la \textbf{Tasa de Impuesto a la Renta}.
\end{itemize}
\end{resultadobox}

\subsection{Nueva Producción de Equilibrio}

Recordamos la condición de equilibrio:
\begin{equation}
DA = Y \tag{6}
\end{equation}

Reemplazamos (5) en (6):
\begin{align}
\barvar{A}_0 + c(1-t)Y &= Y \notag \\
Y - c(1-t)Y &= \barvar{A}_0 \notag \\
Y\left(1 - c(1-t)\right) &= \barvar{A}_0 \notag
\end{align}

\begin{resultadobox}[Nueva Producción de Equilibrio]
\begin{equation}
Y^* = \frac{\barvar{A}_0}{1 - c(1-t)}
\end{equation}
con el \textbf{Nuevo Multiplicador}:
\[
\alpha_m = \frac{1}{1 - c(1-t)}
\]
\end{resultadobox}

\subsection{Comparación: Antes y Después del Impuesto a la Renta}

\begin{center}
\begin{tabular}{lcc}
 & \textbf{Antes} (impuesto suma fija) & \textbf{Ahora} (impuesto a la renta) \\
\hline
$\barvar{A}_0$ & $\barvar{C} - c\barvar{T} + c\barvar{Tr} + \barvar{I} + \barvar{G}$ & $\barvar{C} + c\barvar{Tr} + \barvar{I} + \barvar{G}$ \\[8pt]
Pendiente DA & $c$ & $c(1-t)$ \\[8pt]
Multiplicador & $\dfrac{1}{1-c}$ & $\dfrac{1}{1-c(1-t)}$ \\[8pt]
\end{tabular}
\end{center}

\begin{center}
\begin{tikzpicture}
    \begin{axis}[
        axis lines=middle,
        xlabel={$Y$}, ylabel={$DA$},
        xmin=-1, xmax=14, ymin=-1, ymax=14,
        xtick=\empty, ytick=\empty,
        every axis x label/.style={at={(ticklabel* cs:1)}, anchor=west},
        every axis y label/.style={at={(ticklabel* cs:1)}, anchor=south},
        width=12cm, height=9cm,
    ]
    % Línea de 45 grados
    \addplot[domain=0:13, thick, blue] {x};
    \node[above left, blue] at (axis cs:12,12) {$DA = Y$};

    % DA antes (pendiente c = 0.6, Ao mayor porque incluye -cT)
    \addplot[domain=0:13, thick, green!60!black] {3 + 0.7*x};
    \node[right, green!60!black] at (axis cs:10,10) {$DA = \barvar{A}_0 + cY$ (Antes)};

    % DA ahora (pendiente c(1-t) = 0.6*0.7 = 0.42, Ao sin cT)
    \addplot[domain=0:13, thick, red!70!black] {2.5 + 0.5*x};
    \node[right, red!70!black] at (axis cs:10.5,7.7) {$DA = \barvar{A}_0 + c(1\!-\!t)Y$ (Ahora)};

    % Equilibrios
    \addplot[only marks, mark=*, mark size=3pt, green!50!black] coordinates {(10, 10)};
    \addplot[only marks, mark=*, mark size=3pt, red!60!black] coordinates {(5, 5)};

    % Líneas punteadas
    \addplot[dashed, gray] coordinates {(10,0) (10,10)};
    \addplot[dashed, gray] coordinates {(5,0) (5,5)};

    % Etiquetas
    \node[left] at (axis cs:0,3) {$\barvar{A}_0$};
    \node[left] at (axis cs:0,2.5) {$\barvar{A}_0'$};
    \node[below, red!60!black] at (axis cs:5,0) {$Y^*{}'$};
    \node[below, green!50!black] at (axis cs:10,0) {$Y^*$};

    \end{axis}
\end{tikzpicture}
\end{center}

\begin{explicacionbox}[Conclusión]
El impuesto a la renta \textbf{aminora el multiplicador}. La pendiente de la función de demanda agregada pasa de $c$ a $c(1-t) < c$, lo que reduce el efecto multiplicador sobre la producción de equilibrio. El multiplicador pasa de $\frac{1}{1-c}$ a $\frac{1}{1-c(1-t)} < \frac{1}{1-c}$.
\end{explicacionbox}

\end{document}
